\documentclass[11pt]{amsart}

\usepackage[T1]{fontenc}
\usepackage[utf8]{inputenc}
\usepackage{amsmath,amssymb,amsthm}
\usepackage{booktabs}
\usepackage[a4paper,left=28mm,right=28mm,top=27mm,bottom=27mm]{geometry}
\usepackage[colorlinks=true,linkcolor=blue,citecolor=blue,urlcolor=blue]{hyperref}
\usepackage{microtype}

\newtheorem{theorem}{Theorem}[section]
\newtheorem{lemma}[theorem]{Lemma}
\newtheorem{corollary}[theorem]{Corollary}
\theoremstyle{definition}
\newtheorem{assumption}[theorem]{Assumption}
\newtheorem{definition}[theorem]{Definition}
\theoremstyle{remark}
\newtheorem{remark}[theorem]{Remark}
\newtheorem{example}[theorem]{Example}

\DeclareMathOperator{\supp}{supp}
\DeclareMathOperator{\sign}{sign}
\DeclareMathOperator{\atom}{atom}
\DeclareMathOperator{\cont}{cont}

\newcommand{\R}{\mathbb{R}}
\newcommand{\M}{\mathcal{M}}
\newcommand{\N}{\mathcal{N}}
\newcommand{\TV}{\mathrm{TV}}
\newcommand{\weakn}{\rightharpoonup_{\mathrm n}}

\title{A Narrow-Convergence Existence Proof with Parameter-Moment Regularization}
\author{}
\date{\today}

\begin{document}

\begin{abstract}
This note formulates and proves an existence theorem for a modified version of
the unbounded-parameter sparse-neural-network problem from the thesis.  The
modification is the addition of a coercive parameter-moment term
\[
  \beta\int_{\R^{d+1}}\bigl(1+|\omega|^p\bigr)\,d|\mu|(\omega).
\]
The moment bound supplies the tightness that is missing from the original
location-blind penalty.  Signed Prokhorov compactness then gives narrow
convergence, while the condition that the moment order dominate the growth of
the activation gives uniform integrability of the network atoms.  The note also
records a finite-support consequence under an additional global
$H^1$-growth assumption of the activation function.  The result concerns the modified objective; it is
not a topology-only reproof of existence for the original objective.
\end{abstract}

\maketitle

\section{Problem and notation}

Let $D\subset\R^d$ be a bounded Lipschitz domain, let
$V\in H^1(D)$, and set
\[
  \Omega:=\R^{d+1},\qquad \omega=(a,b)\in\R^d\times\R.
\]
Given a continuous activation profile $\rho:\R\to\R$, define the ridge atom
\[
  K(\omega)(x):=\sigma(x,\omega):=\rho(a\cdot x+b),
  \qquad x\in D.
\]

For $p>0$, introduce the coercive weight
\[
  w_p(\omega):=1+|\omega|^p
\]
and the weighted space of finite signed Radon measures
\[
  \M_p(\Omega)
  :=
  \left\{
    \mu\in\M(\Omega):
    \Psi_p(\mu):=\int_\Omega w_p(\omega)\,d|\mu|(\omega)<\infty
  \right\}.
\]

\begin{definition}[Narrow convergence for signed measures]
We write $\mu_n\weakn\mu$ if
\[
  \int_\Omega f\,d\mu_n\longrightarrow\int_\Omega f\,d\mu
  \qquad\text{for every }f\in C_b(\Omega).
\]
\end{definition}

Suppose for the moment that
\begin{equation}
  \label{eq:value-growth}
  |\rho(z)|\leq C_\rho(1+|z|)^k
  \qquad (z\in\R)
\end{equation}
for some $k\geq0$.  Since $D$ is bounded, there is a constant $C_D$ such that
\begin{equation}
  \label{eq:atom-L2-growth}
  \|K(\omega)\|_{L^2(D)}
  \leq C_D(1+|\omega|)^k .
\end{equation}
The map $K:\Omega\to L^2(D)$ is continuous: if $\omega_n\to\omega$, then all
preactivations $a_n\cdot x+b_n$ remain in one compact interval for
$x\in D$, and uniform continuity of $\rho$ on that interval gives
$K(\omega_n)\to K(\omega)$ uniformly on $D$.
If $p>k$ and $\mu\in\M_p(\Omega)$, the map
$\omega\mapsto K(\omega)$ is Bochner integrable in $L^2(D)$ with respect to
$\mu$, and we may define
\[
  \N\mu:=\int_\Omega K(\omega)\,d\mu(\omega)\in L^2(D).
\]
As in the thesis, the gradient is understood distributionally and the fidelity
is
\[
  L(\mu):=
  \begin{cases}
  \dfrac12\|\N\mu-V\|_{L^2(D)}^2
  +\dfrac12\|\nabla(\N\mu)-\nabla V\|_{L^2(D)}^2,
  &\N\mu\in H^1(D),\\[5pt]
  +\infty,&\N\mu\notin H^1(D).
  \end{cases}
\]

Let $\mu=\mu_{\cont}+\mu_{\atom}$ denote its continuous/atomic
decomposition, where
\[
  \mu_{\atom}
  =
  \sum_{\omega\in\atom\mu}\mu(\{\omega\})\delta_\omega.
\]
For a function $\phi:\R_+\to\R_+$, define the thesis penalty
\[
  \Phi_1(\mu)
  :=
  \|\mu_{\cont}\|_{\M(\Omega)}
  +
  \sum_{\omega\in\atom\mu}
  \phi\bigl(|\mu(\{\omega\})|\bigr).
\]

\begin{assumption}[Penalty assumptions]
\label{ass:phi}
The function $\phi$ satisfies the thesis assumptions:
\begin{enumerate}
\item $\phi(0)=0$, $\phi$ is differentiable on $\R_+$, and
      $\phi'(0)=1$;
\item $\phi$ is concave, nondecreasing, and coercive.
\end{enumerate}
\end{assumption}

For $\alpha,\beta>0$, the modified objective is
\begin{equation}
  \label{eq:modified-objective}
  J_{\beta,p}(\mu)
  :=
  L(\mu)+\alpha\Phi_1(\mu)+\beta\Psi_p(\mu),
  \qquad \mu\in\M_p(\Omega),
\end{equation}
and it is extended by $+\infty$ to
$\M(\Omega)\setminus\M_p(\Omega)$.

\begin{remark}[This is a genuine model modification]
For a finite network
\[
  \mu=\sum_{j=1}^N c_j\delta_{\omega_j},
\]
the new term is
\[
  \beta\Psi_p(\mu)
  =
  \beta\sum_{j=1}^N |c_j|\bigl(1+|\omega_j|^p\bigr).
\]
Thus the cost of a distant neuron is weighted by the magnitude of its output
coefficient.  The term is a weighted total-variation penalty, or equivalently
a parameter-moment penalty.  It is not generated automatically by the narrow
topology.
The choice $1+|\omega|^p$, rather than $(1+|\omega|)^p$, is immaterial for
compactness but is smoother at the origin when $p>1$ and agrees with the
standard polynomial moment used in the cited literature.
\end{remark}

\section{Main existence theorem}

\begin{theorem}[Existence by the narrow direct method]
\label{thm:narrow-existence}
Assume that $D\subset\R^d$ is bounded and Lipschitz,
$V\in H^1(D)$, Assumption~\ref{ass:phi} holds, and
$\rho\in C(\R)$ satisfies the polynomial growth bound
\eqref{eq:value-growth} for some $k\geq0$.  If
\[
  p>k,\qquad \alpha>0,\qquad \beta>0,
\]
then the modified problem
\[
  \inf_{\mu\in\M_p(\Omega)}J_{\beta,p}(\mu)
\]
admits at least one global minimizer.
\end{theorem}

The proof uses three elementary consequences of the moment bound.

\subsection{Moment bounds give narrow compactness}

\begin{lemma}[Tightness and signed Prokhorov compactness]
\label{lem:moment-compactness}
Let $(\mu_n)\subset\M_p(\Omega)$ satisfy
\[
  \sup_n\Psi_p(\mu_n)\leq M<\infty.
\]
Then $(\mu_n)$ has a narrowly convergent subsequence.
\end{lemma}

\begin{proof}
Since $w_p\geq1$,
\[
  \|\mu_n\|_{\M(\Omega)}
  =|\mu_n|(\Omega)
  \leq\Psi_p(\mu_n)
  \leq M.
\]
Moreover, for every $R>0$,
\begin{equation}
  \label{eq:tightness-estimate}
  |\mu_n|(\Omega\setminus B_R)
  \leq
  \frac{1}{1+R^p}
  \int_{\Omega\setminus B_R}w_p\,d|\mu_n|
  \leq\frac{M}{1+R^p}.
\end{equation}
Hence the variations $|\mu_n|$ are uniformly tight.

Write the Jordan decomposition as
$\mu_n=\mu_n^+-\mu_n^-$.  Both nonnegative families
$(\mu_n^+)$ and $(\mu_n^-)$ are uniformly bounded in mass and uniformly
tight because $\mu_n^\pm\leq|\mu_n|$.  Prokhorov's theorem, applied twice and
with a diagonal subsequence, gives nonnegative finite measures
$\lambda^+,\lambda^-$ such that
\[
  \mu_n^+\weakn\lambda^+,
  \qquad
  \mu_n^-\weakn\lambda^-.
\]
It follows that
\[
  \mu_n\weakn\lambda^+-\lambda^-.
\]
The pair $(\lambda^+,\lambda^-)$ need not be the Jordan decomposition of the
limit, but this is irrelevant for compactness.
\end{proof}

\subsection{Lower semicontinuity of the moment}

\begin{lemma}[Narrow lower semicontinuity of $\Psi_p$]
\label{lem:moment-lsc}
If $\mu_n\weakn\mu$, then
\[
  \Psi_p(\mu)
  \leq
  \liminf_{n\to\infty}\Psi_p(\mu_n).
\]
\end{lemma}

\begin{proof}
For $m\in\mathbb{N}$ set
\[
  g_m(\omega):=\min\{w_p(\omega),m\}.
\]
Then $g_m\in C_b(\Omega)$ and $g_m\geq0$.  For every finite signed measure
$\nu$,
\[
  \int_\Omega g_m\,d|\nu|
  =
  \|g_m\nu\|_{\M(\Omega)}
  =
  \sup_{\substack{f\in C_0(\Omega)\\\|f\|_\infty\leq1}}
  \int_\Omega f g_m\,d\nu.
\]
If $\mu_n\weakn\mu$, then
$\int f g_m\,d\mu_n\to\int f g_m\,d\mu$ for every
$f\in C_0(\Omega)$, because $fg_m\in C_0(\Omega)\subset C_b(\Omega)$.
Taking the supremum gives
\[
  \int_\Omega g_m\,d|\mu|
  \leq
  \liminf_{n\to\infty}\int_\Omega g_m\,d|\mu_n|
  \leq
  \liminf_{n\to\infty}\Psi_p(\mu_n).
\]
Finally, $g_m\uparrow w_p$ pointwise, so monotone convergence and then the
supremum over $m$ yield the claim.
\end{proof}

\subsection{Continuity of the value operator}

\begin{lemma}[Uniform moments imply strong value convergence]
\label{lem:value-continuity}
Assume \eqref{eq:atom-L2-growth}, let $p>k$, and suppose
\[
  \mu_n\weakn\mu,
  \qquad
  \sup_n\Psi_p(\mu_n)\leq M.
\]
Then
\[
  \N\mu_n\longrightarrow\N\mu
  \qquad\text{strongly in }L^2(D).
\]
\end{lemma}

\begin{proof}
Lemma~\ref{lem:moment-lsc} gives $\Psi_p(\mu)\leq M$.
Choose $\chi_R\in C_c(\Omega)$ with
\[
  0\leq\chi_R\leq1,\qquad
  \chi_R=1\text{ on }B_R,\qquad
  \supp\chi_R\subset B_{R+1}.
\]
The truncated map
\[
  K_R:\Omega\to L^2(D),\qquad
  K_R(\omega):=\chi_R(\omega)K(\omega),
\]
is continuous and compactly supported.  We first claim that
\begin{equation}
  \label{eq:compact-vector-integral}
  \int_\Omega K_R\,d\mu_n
  \longrightarrow
  \int_\Omega K_R\,d\mu
  \qquad\text{strongly in }L^2(D).
\end{equation}
Indeed, a continuous compactly supported Banach-valued map can be approximated
uniformly by finite-rank maps of the form
\[
  F_\varepsilon(\omega)
  =
  \sum_{j=1}^{J_\varepsilon}f_j(\omega)h_j,
  \qquad
  f_j\in C_c(\Omega),\quad h_j\in L^2(D).
\]
This follows by covering the compact image of $K_R$ with finitely many
$L^2$-balls and using a continuous partition of unity; see
Ryan~\cite[Section~3.2]{ryan-tensor}.  Narrow convergence
gives
\[
  \int_\Omega F_\varepsilon\,d\mu_n
  \longrightarrow
  \int_\Omega F_\varepsilon\,d\mu
  \quad\text{in }L^2(D),
\]
because only finitely many scalar integrals are involved.  The approximation
error is bounded by
\[
  \left\|
    \int_\Omega(K_R-F_\varepsilon)\,d\mu_n
  \right\|_{L^2(D)}
  \leq
  \|K_R-F_\varepsilon\|_\infty\|\mu_n\|_{\M(\Omega)},
\]
uniformly in $n$, and similarly for $\mu$.  This proves
\eqref{eq:compact-vector-integral}.

It remains to remove the cutoff.  For $R\geq1$, $|\omega|>R$, and $p>k$,
\[
  (1+|\omega|)^k
  \leq
  2^k R^{k-p}\bigl(1+|\omega|^p\bigr).
\]
Consequently,
\begin{align*}
  \left\|
    \int_\Omega(1-\chi_R)K\,d\mu_n
  \right\|_{L^2(D)}
  &\leq
  \int_{\Omega\setminus B_R}\|K(\omega)\|_{L^2(D)}\,d|\mu_n|(\omega)\\
  &\leq
  2^k C_D R^{k-p}\Psi_p(\mu_n)\\
  &\leq
  2^k C_D M R^{k-p}.
\end{align*}
The same estimate holds for $\mu$.  First let $n\to\infty$ in the truncated
integrals and then let $R\to\infty$.  Since $k-p<0$, the tail estimate tends
to zero and the claimed strong convergence follows.
\end{proof}

\subsection{Inherited lower semicontinuity of the nonconvex penalty}

\begin{lemma}[Lower semicontinuity of the thesis penalty]
\label{lem:phi-lsc}
Under Assumption~\ref{ass:phi}, if $\mu_n\weakn\mu$, then
\[
  \Phi_1(\mu)
  \leq
  \liminf_{n\to\infty}\Phi_1(\mu_n).
\]
\end{lemma}

\begin{proof}
Narrow convergence implies weak-$*$ convergence against $C_0(\Omega)$.
The sequential weak-$*$ lower semicontinuity of $\Phi_1$ under
Assumption~\ref{ass:phi} is the lower-semicontinuity lemma already proved in
the thesis, based on the nonconvex measure-functional theory of
Bouchitt\'e and Buttazzo~\cite{bouchitte-buttazzo}.  Hence that result applies
without modification.
\end{proof}

\subsection{Proof of Theorem~\ref{thm:narrow-existence}}

\begin{proof}[Proof of Theorem~\ref{thm:narrow-existence}]
Since $V\in H^1(D)$,
\[
  J_{\beta,p}(0)=\frac12\|V\|_{H^1(D)}^2<\infty.
\]
Choose a minimizing sequence $(\mu_n)\subset\M_p(\Omega)$ such that
$J_{\beta,p}(\mu_n)\leq C$.  All three terms in
\eqref{eq:modified-objective} are nonnegative, and therefore
\[
  \Psi_p(\mu_n)\leq\frac{C}{\beta}.
\]
Lemma~\ref{lem:moment-compactness} gives, after extracting a subsequence,
\[
  \mu_n\weakn\bar\mu.
\]
By Lemma~\ref{lem:moment-lsc},
$\bar\mu\in\M_p(\Omega)$ and
\[
  \Psi_p(\bar\mu)
  \leq
  \liminf_n\Psi_p(\mu_n).
\]
Lemma~\ref{lem:value-continuity} gives
\begin{equation}
  \label{eq:value-strong}
  \N\mu_n\longrightarrow\N\bar\mu
  \qquad\text{strongly in }L^2(D).
\end{equation}

The bound on the fidelity gives
\[
  \sup_n\|\nabla(\N\mu_n)\|_{L^2(D)}
  \leq
  \|\nabla V\|_{L^2(D)}+\sqrt{2C}.
\]
After another extraction,
\[
  \nabla(\N\mu_n)\rightharpoonup g
  \qquad\text{weakly in }L^2(D)^d
\]
for some $g$.  For every $\psi\in C_c^\infty(D)^d$,
\begin{align*}
  \int_D g\cdot\psi\,dx
  &=
  \lim_{n\to\infty}
  \int_D\nabla(\N\mu_n)\cdot\psi\,dx\\
  &=
  -\lim_{n\to\infty}
  \int_D\N\mu_n\,\operatorname{div}\psi\,dx\\
  &=
  -\int_D\N\bar\mu\,\operatorname{div}\psi\,dx,
\end{align*}
where the last step follows from \eqref{eq:value-strong}.  Hence
$g=\nabla(\N\bar\mu)$ distributionally and
$\N\bar\mu\in H^1(D)$.

The value part of the fidelity is continuous by
\eqref{eq:value-strong}, while the gradient part is weakly lower
semicontinuous.  Therefore
\[
  L(\bar\mu)\leq\liminf_{n\to\infty}L(\mu_n).
\]
Lemma~\ref{lem:phi-lsc} and Lemma~\ref{lem:moment-lsc} give the same
inequality for the two regularizers.  Combining the three terms,
\[
  J_{\beta,p}(\bar\mu)
  \leq
  \liminf_{n\to\infty}J_{\beta,p}(\mu_n)
  =
  \inf_{\mu\in\M_p(\Omega)}J_{\beta,p}(\mu).
\]
Thus $\bar\mu$ is a global minimizer.
\end{proof}

\section{A finite-support consequence}

The existence theorem only uses growth of the value atoms in $L^2(D)$.  To
make the thesis's finite-support argument unconditional, the moment must
instead dominate the growth of the full $H^1$ atom.

\begin{assumption}[$H^1$ atom growth]
\label{ass:H1-growth}
The map $K:\Omega\to H^1(D)$ is continuous and, for some $q\geq0$,
\begin{equation}
  \label{eq:H1-growth}
  \|K(\omega)\|_{H^1(D)}
  \leq C_K(1+|\omega|)^q
  \qquad(\omega\in\Omega).
\end{equation}
\end{assumption}

\begin{assumption}[Strict local concavity]
\label{ass:strict-concavity}
There exist $\gamma>0$ and $z_0>0$ such that
\[
  \phi'(z_2)-\phi'(z_1)
  \leq
  -\gamma(z_2-z_1)
  \qquad
  (0\leq z_1\leq z_2\leq z_0).
\]
\end{assumption}

\begin{theorem}[Finite support without a separate no-escape condition]
\label{thm:finite-support}
Assume Assumptions~\ref{ass:phi}, \ref{ass:H1-growth}, and
\ref{ass:strict-concavity}.  Let $p>q$, and let
$\bar\mu\in\M_p(\Omega)$ be a local minimizer of $J_{\beta,p}$ with respect
to the total-variation norm, with $J_{\beta,p}(\bar\mu)<\infty$.
Then $\bar\mu$ is finitely supported.
\end{theorem}

\begin{proof}
Set
\[
  r:=\N\bar\mu-V\in H^1(D)
\]
and define the dual field
\[
  P(\omega)
  :=
  \langle r,K(\omega)\rangle_{H^1(D)}.
\]
Assumption~\ref{ass:H1-growth} implies that $P$ is continuous and
\begin{equation}
  \label{eq:P-growth}
  |P(\omega)|
  \leq
  C_K\|r\|_{H^1(D)}(1+|\omega|)^q.
\end{equation}

\medskip
\noindent
\emph{Step 1: first-order conditions.}
Coefficient perturbations by $\pm t\delta_\omega$ give
\begin{equation}
  \label{eq:weighted-insertion}
  |P(\omega)|
  \leq
  \alpha+\beta w_p(\omega)
  \qquad(\omega\notin\atom\bar\mu).
\end{equation}
If $c=\bar\mu(\{\omega\})\neq0$, perturbing the existing coefficient in both
directions gives
\begin{equation}
  \label{eq:weighted-atom-condition}
  P(\omega)
  =
  -\left[
    \alpha\phi'(|c|)+\beta w_p(\omega)
  \right]\sign(c).
\end{equation}
Finally, applying the two scaling directions to compact restrictions of
$\bar\mu_{\cont}$ and then exhausting $\Omega$ gives
\begin{equation}
  \label{eq:weighted-cont-condition}
  P(\omega)
  =
  -\left[\alpha+\beta w_p(\omega)\right]
  \sign(\bar\mu_{\cont})(\omega)
  \quad\text{for }|\bar\mu_{\cont}|\text{-a.e. }\omega.
\end{equation}
These calculations are the thesis's directional-derivative proof with the
additional derivative
\[
  \frac{d}{dt}\bigg|_{t=0}
  \Psi_p(\bar\mu+t u)
\]
included.  All directions used above are compactly supported, so their
networks lie in $H^1(D)$ by Assumption~\ref{ass:H1-growth}.

\medskip
\noindent
\emph{Step 2: the support is automatically confined.}
Equations \eqref{eq:weighted-atom-condition} and
\eqref{eq:weighted-cont-condition} show that every atom, and
$|\bar\mu_{\cont}|$-almost every point, belongs to
\[
  E
  :=
  \{\omega\in\Omega:\beta w_p(\omega)\leq|P(\omega)|\}.
\]
The set $E$ is closed.  By \eqref{eq:P-growth} and $p>q$, it is also bounded,
and therefore compact.  Hence
\[
  \supp\bar\mu\subset E.
\]
This is the step that replaces the thesis's separate no-escape
condition.

\medskip
\noindent
\emph{Step 3: the continuous part vanishes.}
Suppose, for contradiction, that the positive continuous part
$\bar\mu_{\cont}^+$ is nonzero; the negative case is analogous.  By
\eqref{eq:weighted-cont-condition}, choose
\[
  \widehat\omega\in\supp\bar\mu_{\cont}^+,
  \qquad
  \bar\mu(\{\widehat\omega\})=0,
  \qquad
  P(\widehat\omega)+\beta w_p(\widehat\omega)=-\alpha.
\]
For $\delta>0$, let
\[
  \nu_\delta
  :=
  \bar\mu_{\cont}^+|_{B_\delta(\widehat\omega)},
  \qquad
  m_\delta:=\nu_\delta(\Omega),
\]
and merge this mass onto its center:
\[
  \mu_\delta
  :=
  \bar\mu-\nu_\delta+m_\delta\delta_{\widehat\omega}.
\]
Then $m_\delta>0$, $m_\delta\to0$, and
$\|\mu_\delta-\bar\mu\|_{\M(\Omega)}\leq2m_\delta\to0$.

Because the support lies in a compact set and
$K:\Omega\to H^1(D)$ is continuous, there is a modulus
$\eta(\delta)\to0$ such that
\[
  \|\N(\mu_\delta-\bar\mu)\|_{H^1(D)}
  \leq
  \eta(\delta)m_\delta.
\]
The quadratic fidelity has the exact expansion
\begin{align*}
  L(\mu_\delta)-L(\bar\mu)
  &=
  \int_\Omega P\,d(\mu_\delta-\bar\mu)
  +
  \frac12
  \|\N(\mu_\delta-\bar\mu)\|_{H^1(D)}^2.
\end{align*}
The linear fidelity term and the change of the moment cancel together:
\begin{align*}
  &\int_\Omega P\,d(\mu_\delta-\bar\mu)
  +\beta\bigl[\Psi_p(\mu_\delta)-\Psi_p(\bar\mu)\bigr]\\
  &\qquad
  =
  m_\delta\bigl[P(\widehat\omega)+\beta w_p(\widehat\omega)\bigr]
  -
  \int_{B_\delta(\widehat\omega)}
  \bigl[P+\beta w_p\bigr]\,d\nu_\delta
  =0.
\end{align*}
The moved mass was continuous and creates one new atom, so
\[
  \Phi_1(\mu_\delta)-\Phi_1(\bar\mu)
  =
  \phi(m_\delta)-m_\delta.
\]
Assumption~\ref{ass:strict-concavity} and $\phi'(0)=1$ imply, for sufficiently
small $m_\delta$,
\[
  \phi(m_\delta)-m_\delta
  \leq
  -\frac{\gamma}{2}m_\delta^2.
\]
Consequently,
\[
  J_{\beta,p}(\mu_\delta)-J_{\beta,p}(\bar\mu)
  \leq
  \frac12\left[\eta(\delta)^2-\alpha\gamma\right]m_\delta^2<0
\]
for small $\delta$, contradicting local minimality.  Thus
$\bar\mu_{\cont}=0$.

\medskip
\noindent
\emph{Step 4: there are only finitely many atoms.}
Suppose instead that $\bar\mu$ has infinitely many atoms.  Since their support
is contained in the compact set $E$, there are distinct atoms
$\omega_i\to\widehat\omega$ whose coefficients have a common sign.  Reversing
signs if necessary, write
\[
  c_i:=\bar\mu(\{\omega_i\})>0.
\]
Finite total variation gives $c_i\to0$.  By
\eqref{eq:weighted-atom-condition}, continuity of $P$ and $w_p$, and
$\phi'(c_i)\to\phi'(0)=1$,
\[
  P(\widehat\omega)
  =
  -\alpha-\beta w_p(\widehat\omega).
\]
For $N\in\mathbb{N}$ set
\[
  C_N:=\sum_{i=N}^\infty c_i,\qquad
  \delta_N:=\sup_{i\geq N}|\omega_i-\widehat\omega|,
\]
and merge the tail:
\[
  \mu_N
  :=
  \bar\mu-\sum_{i=N}^\infty c_i\delta_{\omega_i}
  +C_N\delta_{\widehat\omega}.
\]
Then $C_N\to0$, $\delta_N\to0$, and
$\|\mu_N-\bar\mu\|_{\M(\Omega)}\leq2C_N\to0$.  Uniform continuity of $K$ on
the compact support gives
\[
  \|\N(\mu_N-\bar\mu)\|_{H^1(D)}
  \leq\eta(\delta_N)C_N.
\]
At a pre-existing atom at $\widehat\omega$, the triangle inequality gives the
upper bounds
\begin{align*}
  \Psi_p(\mu_N)-\Psi_p(\bar\mu)
  &\leq
  -\sum_{i=N}^\infty c_i w_p(\omega_i)
  +C_N w_p(\widehat\omega),\\
  \Phi_1(\mu_N)-\Phi_1(\bar\mu)
  &\leq
  -\sum_{i=N}^\infty\phi(c_i)+\phi(C_N).
\end{align*}
Combining these estimates with the exact quadratic expansion of $L$ and the
weighted atom condition gives
\begin{align*}
  J_{\beta,p}(\mu_N)-J_{\beta,p}(\bar\mu)
  &\leq
  -\alpha\sum_{i=N}^\infty
  \bigl[\phi(c_i)-c_i\phi'(c_i)\bigr]\\
  &\quad
  -\alpha\bigl[C_N-\phi(C_N)\bigr]
  +\frac12\eta(\delta_N)^2C_N^2.
\end{align*}
Concavity and $\phi(0)=0$ give
$\phi(c)-c\phi'(c)\geq0$.  Assumption~\ref{ass:strict-concavity} gives
$C_N-\phi(C_N)\geq\frac{\gamma}{2}C_N^2$ for large $N$.  Therefore
\[
  J_{\beta,p}(\mu_N)-J_{\beta,p}(\bar\mu)
  \leq
  \frac12\left[\eta(\delta_N)^2-\alpha\gamma\right]C_N^2<0
\]
for large $N$, again contradicting local minimality.  The atomic support is
therefore finite.
\end{proof}

\begin{corollary}[Existence of a finite network]
\label{cor:finite-global-minimizer}
Assume the hypotheses of Theorem~\ref{thm:narrow-existence},
Assumptions~\ref{ass:H1-growth} and \ref{ass:strict-concavity}, and
\[
  p>\max\{k,q\}.
\]
Then the modified problem admits a finitely supported global minimizer.  In
particular, the existence and finite-support conclusions compose without a
separate no-escape condition.
\end{corollary}

\begin{proof}
Theorem~\ref{thm:narrow-existence} supplies a global minimizer
$\bar\mu\in\M_p(\Omega)$ with finite objective.  A global minimizer is a
total-variation local minimizer, and $p>q$ allows
Theorem~\ref{thm:finite-support} to be applied to $\bar\mu$.
\end{proof}

\section{Activation exponents}

The following safe exponents follow from elementary pointwise bounds on a
bounded domain.  The exponent $k$ controls the value atom in $L^2(D)$ and is
enough for Theorem~\ref{thm:narrow-existence}; the exponent $q$ controls the
full atom in $H^1(D)$ and is needed for
Theorem~\ref{thm:finite-support}.
\[
\begin{array}{lccc}
\toprule
\text{activation} & k & \text{safe }q
& \text{one choice giving existence and finite support}\\
\midrule
\tanh z              &0&1&p>1\\
e^{-z^2}             &0&1&p>1\\
\log(1+e^z)          &1&1&p>1\\
(z_+)^m              &m&m&p>m\\
\bottomrule
\end{array}
\]
For bounded activations such as $\tanh$ and the Gaussian, existence alone
requires only $p>0$.  The safer condition $p>1$ is shown in the table because
it also dominates the elementary global $H^1$ growth estimate.

\section{What fails when the moment term is absent}

\begin{example}[Loss of narrow compactness]
Let $\omega_n=(0,n)$ and $\mu_n=\delta_{\omega_n}$.  Then
\[
  \|\mu_n\|_{\M(\Omega)}=1,
  \qquad
  \Phi_1(\mu_n)=\phi(1),
\]
so the original regularizer is uniformly bounded.  Nevertheless,
$(\mu_n)$ has no narrowly convergent subsequence.  Indeed, compactly supported
test functions force every putative limit to be zero, while the bounded
continuous test function $1$ preserves the mass:
\[
  \int_\Omega1\,d\mu_n=1.
\]
Thus bounded mass without a location-coercive term does not give narrow
compactness.
\end{example}

\begin{example}[Narrow convergence does not control an unbounded dictionary]
For softplus, let
\[
  \rho(z)=\log(1+e^z),\qquad
  \omega_n=(0,n),\qquad
  \mu_n=\frac{1}{\rho(n)}\delta_{\omega_n}.
\]
Then $\mu_n\weakn0$, because
$\|\mu_n\|_{\M(\Omega)}=\rho(n)^{-1}\to0$.  However,
\[
  \N\mu_n\equiv1
  \qquad\text{while}\qquad
  \N0=0.
\]
For $V\equiv1$, the original objective satisfies
\[
  L(\mu_n)=0,\qquad
  \Phi_1(\mu_n)=\phi(\rho(n)^{-1})\longrightarrow0,
\]
which yields the nonattainment mechanism recorded in the thesis.  With the
moment term,
\[
  \Psi_p(\mu_n)
  =
  \frac{1+n^p}{\rho(n)}
  \asymp n^{p-1},
\]
so every $p>1$ prices out this escape.
\end{example}

\section{Relation to the literature}

The compactness mechanism is classical: Prokhorov's theorem converts uniform
tightness and a total-mass bound into narrow precompactness.  A coercive moment
bound supplies tightness through the elementary estimate
\eqref{eq:tightness-estimate}; see, for example,
Bogachev~\cite{bogachev-measure}.

In the neural-network literature, Bartolucci, Carioni, Iglesias, Korolev,
Naldi, and Vigogna~\cite{bartolucci-et-al} study signed measures on unbounded
parameter spaces and explicitly regularize by total variation and parameter
moments.  Their Theorem~3.7 proves compactness in the stronger unbalanced
Kantorovich--Rubinstein norm for moments of order strictly larger than one.
The present proof only seeks narrow compactness, so the compactness step itself
needs any coercive moment.  The additional restriction $p>k$ comes instead
from uniform integrability of a dictionary growing like $|\omega|^k$.

\begin{thebibliography}{9}

\bibitem{bartolucci-et-al}
F.~Bartolucci, M.~Carioni, J.~A.~Iglesias, Y.~Korolev, E.~Naldi, and
S.~Vigogna,
\emph{A Lipschitz spaces view of infinitely wide shallow neural networks},
SIAM J. Math. Anal. \textbf{58} (2026), no.~3, 2786--2828.
\href{https://doi.org/10.1137/24M1718615}{doi:10.1137/24M1718615}.

\bibitem{bogachev-measure}
V.~I.~Bogachev,
\emph{Measure Theory, Vol.~II},
Springer, Berlin, 2007.

\bibitem{bouchitte-buttazzo}
G.~Bouchitt\'e and G.~Buttazzo,
\emph{New lower semicontinuity results for nonconvex functionals defined on
measures},
Nonlinear Anal. \textbf{15} (1990), no.~7, 679--692.

\bibitem{ryan-tensor}
R.~A.~Ryan,
\emph{Introduction to Tensor Products of Banach Spaces},
Springer Monographs in Mathematics,
Springer-Verlag, London, 2002, Section~3.2, pp.~49--50.
\href{https://doi.org/10.1007/978-1-4471-3903-4}
{doi:10.1007/978-1-4471-3903-4}.

\end{thebibliography}

\end{document}
